% S04-M2 (package abhijit/s04-math): the Lambda-adic Brandt module, freeness, Hida's lemma, control (a), factor 2.
% Labels: def:s04-Mord, prop:s04-freeA, lem:s04-hida, thm:s04-control, rem:s04-hida86.
% thm:s04-control(b) is stated here (one theorem, one label); its dictionary is S04-M4 (Remark rem:s04-controlb in hsieh.tex).
\section{Hida theory on the definite side}

Notation as in \S\ref{ss:s04-notation}. Everything in this section holds verbatim over $\cO=\cOE$ in place of $\Zp$, because
$\cO[X]=\Zp[X]\otimes_{\Zp}\cO$, $\cO$ is $\Zp$-free, and the ordinary projector commutes with $\otimes_{\Zp}\cO$ (it is the limit of
$U_p^{n!}$, and $U_p^{n!}\otimes1$ converges to $e\otimes1$).

\begin{definition}[$\Lambda$-adic Brandt module; seed \S7]\label{def:s04-Mord}
For $U$ of level prime to $p$: $\Mord(U):=\varprojlim_{s\ge1}e\,\cOE[X_s(U)]$, limit along $\pi_*$ (which commutes with $U_p$ by
Lemma~\ref{lem:s01-Up}(b), and with all $T_\ell,S_\ell,U_\ell,\Pi_\ell,\langle d\rangle$). $\Lam$ acts through the adelic diamonds:
$[u]\cdot a:=\langle u\rangle a$, $u\in1+p\Zp$. The finite group $\mu_{p-1}$ acts as well, and $\Mord(U)$ decomposes integrally into
$\mu_{p-1}$-eigenpieces $\Mord(U)^{(\chi)}$. For $U=\Uone M$ also $(\Z/M)^\times$ acts; if $p\nmid\varphi(M)$, $\Mord(U)$ decomposes
integrally into eigenpieces for $\mu_{p-1}\times(\Z/M)^\times$, and if $p\mid\varphi(M)$ only after inverting $p$ (as in
Thm.~\ref{thm:s04-control}(b)).
\end{definition}
\noindent\emph{Well-definedness.} $\pi_*$ commutes with $U_p$, hence with $e$ (functoriality in Def.~\ref{def:s04-e}), so
$(e\,\cO[X_s])_{s\ge1}$ is an inverse system. $\delta_u=\iota_p^{-1}\mathrm{diag}(1,u)$ normalises every $UU_s$, and $\pi_*[x\delta_u]=[x\delta_u]$,
so the diamonds commute with $\pi_*$; on level $s$ the group $1+p\Zp$ acts through $\Delta_s$, i.e.\ $e\,\cO[X_s]$ is an
$\cO[\Delta_s]$-module, and the transition maps are $\cO[\Delta_{s+1}]$-linear. So $\Mord(U)$ is a module over
$\varprojlim_s\cO[\Delta_s]=\Lam$. The eigenpiece decomposition for a finite abelian group $G$ acting $\cO$-linearly needs $\cO\supseteq\mu_{|G|}$ and
$|G|\in\cO^\times$, and is then given by the usual idempotents $|G|^{-1}\sum_g\chi^{-1}(g)g$. For $G=\mu_{p-1}$, $p\nmid p-1$; for
$G=\mu_{p-1}\times(\Z/M)^\times$, $|G|=(p-1)\varphi(M)$, which is prime to $p$ iff $p\nmid\varphi(M)$ ($p\nmid M$ does not suffice:
$M=11$, $p=5$). Otherwise the idempotents exist only in $\cO[1/p][G]$. \stproved.
\emph{Lean:} \stformal{} for $E=\Q_p$ and $U=\Uone M$, as compatible families (D-lim, D-$\Lam$): \lean{Langlandscape.Ordinary.Mord} (with levels
\lean{Langlandscape.Ordinary.TowerHyp.MordLevel} and transitions \lean{Langlandscape.Ordinary.TowerHyp.mordTrans}), the ring $\Lam$
as \lean{Langlandscape.Ordinary.Lam}, and the $\Lam$-module structure \lean{Langlandscape.Ordinary.Mord.lamModule}, in which the system of
diamond representations $\rho$ commuting with $\pi_*$ is an explicit argument (the well-definedness paragraph above is thus a hypothesis there, not
a Lean theorem). Not formalized: general $\cO_E$ and the $\mu_{p-1}\times(\Z/M)^\times$ eigenpiece remark.

\subsection{Freeness at finite level}

\begin{proposition}[freeness at finite level; seed Prop.~free (a)]\label{prop:s04-freeA}
For $s\ge1$, $e\,\Zp[X_s(U)]$ is a free $\Zp[\Delta_s]$-module. \stproved{} (seed): $\Zp[X_s]$ is $\Zp[\Delta_s]$-free since
$\Delta_s$ acts freely; $U_p$ commutes with $\Delta_s$ (Lemma~\ref{lem:s01-Up}(e)), so $e\Zp[X_s]$ is a direct summand, hence
projective, hence free over the local ring $\Zp[\Delta_s]$. Only $\Delta_s$, not $(\Z/p^s)^\times$, is used.
\end{proposition}
\begin{proof}
\emph{Freeness of $\Zp[X_s]$.} $\Delta_s$ acts freely on $X_s$ (N2); choosing one point $x_i$ in each orbit gives
$\Zp[X_s]=\bigoplus_i\Zp[\Delta_s]\,[x_i]$, a free $\Zp[\Delta_s]$-module.

\emph{$e$ is $\Zp[\Delta_s]$-linear.} $U_p$ commutes with every $\langle d\rangle$ (Lemma~\ref{lem:s01-Up}(e)), so $U_p$, its powers
and their limit $e$ are $\Zp[\Delta_s]$-linear. Hence $\Zp[X_s]=e\Zp[X_s]\oplus(1-e)\Zp[X_s]$ is a decomposition of
$\Zp[\Delta_s]$-modules, and $e\Zp[X_s]$ is a finitely generated projective $\Zp[\Delta_s]$-module.

\emph{$\Zp[\Delta_s]$ is local.} $\Delta_s$ is a finite $p$-group (of order $p^{s-1}$). The augmentation ideal $I_\F$ of $\F_p[\Delta_s]$
is generated by the $g-1$, $g\in\Delta_s$, and $(g-1)^{p^{s-1}}=g^{p^{s-1}}-1=0$ in characteristic $p$; a finitely generated ideal of a
commutative ring generated by nilpotents is nilpotent, so $\F_p[\Delta_s]$ has the unique prime ideal $I_\F$. Since $\Zp[\Delta_s]$
is finite over the complete local ring $\Zp$, every maximal ideal contains $p$, hence corresponds to a maximal ideal of
$\F_p[\Delta_s]$; there is only one, $\m_s=(p)+I$ with $I$ the augmentation ideal.

\emph{Projective implies free.} A finitely generated projective module $P$ over a local ring $R$ is free: lift a basis of
$P/\m P$ to $P$; by Nakayama it generates, giving $0\to K\to R^n\to P\to0$; this splits, so $K$ is a finitely generated direct summand
with $K/\m K=0$ (compare dimensions over $R/\m$), hence $K=0$.
\end{proof}
\noindent\emph{Why only $\Delta_s$.} The prime-to-$p$ part $\mu_{p-1}$ of $(\Z/p^s)^\times$ need not act freely on $X_s$ (stabilisers in
$\mu_{p-1}$ can occur), and $\Zp[\mu_{p-1}]$ is not local; both steps of the proof would fail.
\emph{Lean:} \stformal{} for $\cO=\Zp$, $U=\Uone M$, $p\ge5$: \lean{Langlandscape.Ordinary.ordFree\_finite} (free and finite over $\Zp[\Delta_s]$; its argument
\texttt{hdec : Hyp.PadicOrdDecomp p} is discharged by \lean{Langlandscape.Ordinary.padicOrdDecomp\_holds}, so the statement is unconditional), via
\lean{Langlandscape.Ordinary.isLocalRing\_monoidAlgebra\_padicInt}, \lean{Langlandscape.Ordinary.permMod\_free} and
\lean{Langlandscape.Ordinary.free\_ordPartG}; also over $(\Z/p^m)[\Delta_s]$ as \lean{Langlandscape.Ordinary.ordFree\_finite\_zmod}. Not formalized:
general $\cO_E$ and levels $U\ne\Uone M$.

\subsection{Hida's lemma}

\begin{lemma}[Hida's lemma, definite form]\label{lem:s04-hida}
Let $s\ge1$, $\pi':X'_s\to X_s$ the projection ($[\Kp s:\Kpp s]=p$) and $\pi'':X_{s+1}\to X'_s$ (a $\ker(\Delta_{s+1}\to\Delta_s)$-torsor).
\begin{enumerate}[label=(\alph*)]
\item $\Kpp s\varpi_p\Kpp s=\coprod_j\varpi_j\Kpp s$ with the same $\varpi_j$; so $U_p$ is defined compatibly on $X_{s+1},X'_s,X_s$.
\item For $u\in\Kp s$ and each $j$ there are $i$ and $u_j\in\Kpp s$ with $u\varpi_j=\varpi_iu_j$. Hence, cohomologically,
  $U_p\big(\Zp^{X'_s}\big)\subseteq\pi'^*\Zp^{X_s}$, and homologically $U_p(\ker\pi'_*)=0$.
\item $e\pi'_*:e\Zp[X'_s]\to e\Zp[X_s]$ is an isomorphism; and $\pi''_*$ induces $\Zp[X_{s+1}]\otimes_{\Zp[\Delta_{s+1}]}\Zp[\Delta_s]
  \xrightarrow{\sim}\Zp[X'_s]$.
\item Consequently $\pi_*$ induces $e\Zp[X_{s+1}]\otimes_{\Zp[\Delta_{s+1}]}\Zp[\Delta_s]\xrightarrow{\sim}e\Zp[X_s]$, and
  $\rk_{\Zp[\Delta_s]}e\Zp[X_s]$ is independent of $s\ge1$.
\end{enumerate}
\stproved{} below (all steps elementary). \emph{Lean} (for $\Zp$, $U=\Uone M$, $p\ge5$; \stformal{} unless stated): (a), (b)
\lean{Langlandscape.Local.Kpp\_mul\_varpi\_mul\_Kpp}, \lean{Langlandscape.Local.exists\_varpi\_conj}, and (b) at the level of $\Zp[X]$:
\lean{Langlandscape.Ordinary.Up\_ker\_push\_eq\_zero}; (c) second claim:
\lean{Langlandscape.Local.kppQuotEquiv}, \lean{Langlandscape.Ordinary.piPP\_eq\_iff}, \lean{Langlandscape.Ordinary.piPP\_free}; (d), finite-level
control: \lean{Langlandscape.Ordinary.ord\_control\_finite} ($e\pi_*$ surjective with kernel $e\cap\ker\pi''_*$, i.e.\ the tensor isomorphism in
kernel form, hypothesis \texttt{hdec} discharged by \pathtt{padicOrdDecomp\_holds}). Not formalized: the first claim of (c) as a separate statement,
the tensor product isomorphism as a stated equivalence, and the rank independence of $s$ in (d) (\pathtt{ordRank\_const} does not exist).
\end{lemma}
\begin{proof}
Matrices are at $p$, read through $\iota_p$; all elements below are supported at $p$, so they commute with $U^{(p)}$. Write
$H:=\ker(\Delta_{s+1}\to\Delta_s)=(1+p^s\Zp)/(1+p^{s+1}\Zp)$, of order $p$.

(a) Let $K=\Kpp s$; its elements $\left(\begin{smallmatrix}a&b\\c&d\end{smallmatrix}\right)$ have $c\in p^{s+1}\Zp$, $d\in1+p^s\Zp$,
$a\in\Zp^\times$. As in Lemma~\ref{lem:s01-Up}(a), $K\varpi K/K$ is in bijection with $K/(K\cap\varpi K\varpi^{-1})$, and
$\varpi^{-1}k\varpi=\left(\begin{smallmatrix}a&b/p\\pc&d\end{smallmatrix}\right)$ lies in $K$ iff $b\in p\Zp$ (its lower left entry $pc$
and lower right entry $d$ satisfy the conditions of $K$). For $k\in K$ there is exactly one $j\in\{0,\dots,p-1\}$ with
$b-jd\in p\Zp$, and $n(-j)k\in K\cap\varpi K\varpi^{-1}$, where $n(j)=\left(\begin{smallmatrix}1&j\\0&1\end{smallmatrix}\right)\in K$.
Hence $K\varpi K=\coprod_jn(j)\varpi K=\coprod_j\varpi_jK$; disjointness as in Lemma~\ref{lem:s01-Up}(a). So at the three levels
$UU_{s+1}$, $U^{(p)}\Kpp s$, $UU_s$ (Lemma~\ref{lem:s01-Up}(a) for the outer two) the representatives are the same $\varpi_j$, and
$U_p[x]=\sum_j[x\varpi_j]$ commutes with $\pi'_*$ and $\pi''_*$ (literally the same formula, as in Lemma~\ref{lem:s01-Up}(b)).

(b) Let $u=\left(\begin{smallmatrix}a&b\\c&d\end{smallmatrix}\right)\in\Kp s$, so $c\in p^s\Zp$, $d\in1+p^s\Zp$, $a\in\Zp^\times$.
For $i,j\in\{0,\dots,p-1\}$,
\[\varpi_i^{-1}u\varpi_j=\begin{pmatrix}1/p&-i/p\\0&1\end{pmatrix}\begin{pmatrix}ap&aj+b\\cp&cj+d\end{pmatrix}
=\begin{pmatrix}a-ic&\big(aj+b-i(cj+d)\big)/p\\pc&cj+d\end{pmatrix}.\]
Since $cj+d\equiv d\equiv1\pmod{p^s}$ is a unit, there is exactly one $i$ with $i\equiv(aj+b)(cj+d)^{-1}\pmod p$; for it the matrix
$u_j:=\varpi_i^{-1}u\varpi_j$ is integral, has determinant $\det u\in\Zp^\times$, lower left entry $pc\in p^{s+1}\Zp$ and lower right
entry $cj+d\in1+p^s\Zp$; so $u_j\in\Kpp s$. The map $j\mapsto i=i(j)$ is a bijection: if $i(j)=i(j')$ then
$u\varpi_j\Kpp s=\varpi_i\Kpp s=u\varpi_{j'}\Kpp s$, so $j=j'$ by the disjointness in (a).
\emph{Cohomologically:} for $f\in\Zp^{X'_s}$ (a function on $B^\times\backslash\Bhat^\times$ invariant under $U^{(p)}\Kpp s$) and $u\in\Kp s$,
\[(U_pf)(xu)=\sum_jf(xu\varpi_j)=\sum_jf(x\varpi_{i(j)}u_j)=\sum_jf(x\varpi_{i(j)})=(U_pf)(x),\]
and $U_pf$ is invariant under $U^{(p)}$ as before; so $U_pf$ is invariant under $UU_s$, i.e.\ $U_pf\in\pi'^*\Zp^{X_s}$.
\emph{Homologically:} the homological $U_p$ and $\pi'_*$ are the transposes of the cohomological $U_p$ and $\pi'^*$ for the naive
evaluation pairing (N4), which is perfect. If $\pi'_*a=0$ then for every $f$, writing $U_pf=\pi'^*g$,
$\langle U_pa,f\rangle=\langle a,U_pf\rangle=\langle a,\pi'^*g\rangle=\langle\pi'_*a,g\rangle=0$; so $U_pa=0$. (No weights $w_x$ enter.)

(c) \emph{First claim.} $\pi'$ is surjective, so $0\to\ker\pi'_*\to\Zp[X'_s]\xrightarrow{\pi'_*}\Zp[X_s]\to0$ is exact, and it is
$U_p$-equivariant by (a). $U_p$ is zero on $\ker\pi'_*$ by (b), so $e=\lim U_p^{n!}=0$ there, and exactness of $e$
(Def.~\ref{def:s04-e}) gives $e\pi'_*:e\Zp[X'_s]\xrightarrow{\sim}e\Zp[X_s]$.
\emph{Second claim.} $\Kp{s+1}\subseteq\Kpp s$, and $\Kpp s\to(1+p^s\Zp)/(1+p^{s+1}\Zp)=H$,
$\left(\begin{smallmatrix}a&b\\c&d\end{smallmatrix}\right)\mapsto d$, is a homomorphism (the $(2,2)$ entry of a product is
$cb'+dd'\equiv dd'\bmod p^{s+1}$, as $c\in p^{s+1}\Zp$), surjective (use $\mathrm{diag}(1,d)$) with kernel $\Kp{s+1}$. Hence
$\Kp{s+1}\trianglelefteq\Kpp s$, the fibres of $\pi''$ are the orbits of $\Kpp s/\Kp{s+1}\cong H$ acting by $[x]\mapsto[x\delta_d]$,
$d\in1+p^s\Zp$, and this is the restriction to $H$ of the free $\Delta_{s+1}$-action (N2). So $\pi''$ identifies $X'_s$ with
$X_{s+1}/H$, every fibre being a free $H$-orbit. Therefore
\[\Zp[X_{s+1}]\otimes_{\Zp[\Delta_{s+1}]}\Zp[\Delta_s]=\Zp[X_{s+1}]\otimes_{\Zp[\Delta_{s+1}]}\Zp[\Delta_{s+1}/H]
=\Zp[X_{s+1}]_H=\Zp[X_{s+1}/H]\xrightarrow[\pi''_*]{\sim}\Zp[X'_s],\]
where $(\cdot)_H$ denotes $H$-coinvariants (permutation modules have $\Zp[Y]_H=\Zp[Y/H]$). The isomorphism is $\Delta_s$-equivariant:
$\delta_d$, $d\in1+p\Zp$, normalises $\Kpp s$, and $\delta_d\in\Kpp s$ for $d\in1+p^s\Zp$, so $\Delta_s$ acts on $X'_s$ and $\pi''$ is
equivariant.

(d) $\pi_*=\pi'_*\circ\pi''_*$. $U_p$ commutes with $\langle d\rangle$ on all three levels (Lemma~\ref{lem:s01-Up}(e) for $X_{s+1}$, $X_s$;
for $X'_s$ by the same computation: $\delta_d^{-1}\varpi_j\delta_d=\left(\begin{smallmatrix}p&jd\\0&1\end{smallmatrix}\right)
=\varpi_jn\big(j(d-1)/p\big)$ with $n(\cdot)\in\Kpp s$, and $\delta_d$ normalises $\Kpp s$). For a finitely generated
$\Zp$-module $M$ with a $\Zp[\Delta_{s+1}]$-linear $U_p$, put $\bar M:=M\otimes_{\Zp[\Delta_{s+1}]}\Zp[\Delta_s]=M/IM$,
$I=\ker(\Zp[\Delta_{s+1}]\to\Zp[\Delta_s])$. Then $\bar M=e\bar M\oplus(1-e)\bar M$ with $e\bar M=\overline{eM}$: the decomposition
$M=eM\oplus(1-e)M$ is $\Delta_{s+1}$-stable, so it induces $\bar M=\overline{eM}\oplus\overline{(1-e)M}$ with $U_p$-stable summands. $U_p$ is surjective on
$eM$, hence on its quotient $\overline{eM}$, and a surjective endomorphism of a finitely generated module over a commutative ring is
injective (Vasconcelos); so $U_p$ is invertible on $\overline{eM}$. $U_p$ is topologically nilpotent on $\overline{(1-e)M}$, since
$U_p^n((1-e)M)\subseteq p^kM$ for $n$ large. By uniqueness of the decomposition, $e_{\bar M}$ is the projection onto $\overline{eM}$. Applying this
to $M=\Zp[X_{s+1}]$ and combining with (c):
\[e\Zp[X_{s+1}]\otimes_{\Zp[\Delta_{s+1}]}\Zp[\Delta_s]=e\big(\Zp[X_{s+1}]\otimes_{\Zp[\Delta_{s+1}]}\Zp[\Delta_s]\big)
\xrightarrow[\pi''_*]{\sim}e\Zp[X'_s]\xrightarrow[\pi'_*]{\sim}e\Zp[X_s],\]
and the composite is induced by $\pi_*$. \emph{Rank.} By Prop.~\ref{prop:s04-freeA}, $e\Zp[X_{s+1}]\cong\Zp[\Delta_{s+1}]^{r_{s+1}}$, so
the left side is $\Zp[\Delta_s]^{r_{s+1}}$, and it is isomorphic to $e\Zp[X_s]\cong\Zp[\Delta_s]^{r_s}$; ranks of free modules over the
non-zero commutative ring $\Zp[\Delta_s]$ are well defined, so $r_{s+1}=r_s$.
\end{proof}
\noindent This is the precise answer to the seed's remark after Prop.~free (``fibres have size $p^2$ while
$\ker(\Delta_{s+1}\to\Delta_s)$ has order $p$''): the extra factor $p$ is the step $X'_s\to X_s$, and $e$ collapses it.


\subsection{\texorpdfstring{$\Lambda$}{Lambda}-freeness and control}

\begin{theorem}[$\Lambda$-freeness and control; seed Prop.~free (b); \texttt{if:s04-Mord}]\label{thm:s04-control}
Let $U$ be of level prime to $p$ and $r:=\rk_{\Zp}e\,\Zp[X_1(U)]$.
\begin{enumerate}[label=(\alph*)]
\item $\Mord(U)$ is a free $\Lam$-module of rank $r$, and for every $s\ge1$,
  $\Mord(U)\otimes_\Lam\cOE[\Delta_s]\xrightarrow{\sim}e\,\cOE[X_s(U)]$, compatibly with all Hecke operators.
\item (weights; $U=\Uone M$) For $k\ge2$, $\Mord(U)\otimes_\Lam\Lam/P^{\rm ad}_k$ is dual to $e\,\Brandt{U\Kp1}{\Vk\otimes\cOE}$ (Def.~\ref{def:s04-Pk}).
  On the component where $\mu_{p-1}$ acts by $\omega^{-(k-2)}$ it is dual to $e\,\Brandt{U\cdot\iota_p^{-1}\Uzero{p}_p}{\Vk\otimes\cOE}$, for every $k\ge2$.
  For $k\ge3$ the latter equals $e\,\Brandt{U}{\Vk\otimes\cOE}$ by $p$-stabilisation (\stexp{} in our normalisation).
  The first sentence is claimed integrally (over $\cOE$) only if $p\nmid\varphi(M)$; for $p\mid\varphi(M)$ it holds after inverting $p$,
  and its integral form is \stsketch{} (Remark~\ref{rem:s04-controlb}(iv)). For a general level $U$ prime to $p$, (b) is \stexp{}
  (source: Hida 1988, Thm.~9.4, not checked here).
  \provisional{F4: (b) restated so that the character piece and the level match; before: ``For $k\ge2$ and $\chi$ the component on which $\mu_{p-1}$ acts by $\omega^{-(k-2)}$, $\Mord(U)^{(\chi)}\otimes_\Lam\Lam/P^{\rm ad}_k$ is dual to $e\,M(U\Kp1;\Vk\otimes\cOE)$; for $k\ge3$ one may replace $U\Kp1$ by $U\cdot\iota_p^{-1}\Uzero p_p$ (Hida's control at $\Gamma_0(p)$-level; expected)''; the left side was one $\mu_{p-1}$-eigenpiece and the right side the whole level-$U\Kp1$ module.}
\end{enumerate}
Status: (a) \stproved{} from Lemma~\ref{lem:s04-hida} and Prop.~\ref{prop:s04-freeA} (choose compatible bases by Nakayama;
$\Lam=\varprojlim\cOE[\Delta_s]$); proof below, which uses only stage-01 results and elementary commutative algebra (raised from
\stpmc{} after the chapter audit S04-M9). Also \stcited{} \lit{hsieh\_2017\_hida-families-and-p-adic-triple-product-l-functions\_arxiv}{Thm.~4.2 and
its proof, p.~32: $\varprojlim_\alpha e\cO[X_\alpha]$ is $\Lam[\Delta_p]$-free, via [Hid88b, Cor.~10.1]} (Hsieh's ``Cor.~10.1'' is
\emph{Theorem}~10.1 in Hida's paper; there is no Cor.~10.1) and
\lit{longo-vigni\_2011\_control-theorems-for-quaternionic-hida-families\_arxiv}{Prop.~3.6 and its proof, p.~28: $\varprojlim_r
S_2(U_r,\cO_f)^{\rm ord}$ is $\Lam$-free by [Hid88b, Cor.~10.4]}.
(b) for $U=\Uone M$: \stcited{}, integrally if $p\nmid\varphi(M)$ (otherwise the integral form is \stsketch); general $U$: \stexp.
\lit{hsieh\_2017\_hida-families-and-p-adic-triple-product-l-functions\_arxiv}{Thm.~4.2(2), p.~32 (forms side; via
[Hid88b, Thm.~9.4])}; \lit{longo-vigni\_2011\_control-theorems-for-quaternionic-hida-families\_arxiv}{Thm.~3.7, p.~28 (eigen-components,
over a field)}; the dictionary to our $\NUp$, $\Vk$ and $P^{\rm ad}_k$ is Remark~\ref{rem:s04-controlb} (S04-M4), whose key
computation is Def.~\ref{def:s04-Pk}. The original source is
\lit{hida\_1988\_p-adic-hecke-algebras-for-gl2-over-totally-real-fields}{Thm.~8.6, p.~362 (weight independence of $\mathcal V^{\rm ord}$);
Thm.~9.4, p.~369 (control); Thm.~10.1, p.~370 (freeness over $\cO[[Z_a]]$ under $p$-divisibility); Cor.~10.4, p.~373 ($\Lam$-freeness,
$p>2$)}: the theorem headers and page numbers are located in the OCR text, but the statements are only partly legible, so their exact
hypotheses are taken from Hsieh and LV11.
\emph{Lean status of (a):} the finite-level part is \stformal{} for $\Zp$, $U=\Uone M$, $p\ge5$: control as surjectivity plus kernel
(\lean{Langlandscape.Ordinary.ord\_control\_finite}, \lean{Langlandscape.Ordinary.TowerHyp.mordTrans\_surjective},
\lean{Langlandscape.Ordinary.Mord.proj\_surjective}), freeness at level $s$ (\lean{Langlandscape.Ordinary.ordFree\_finite}). \emph{Not} formalized:
the constancy of the rank in $s$ (\pathtt{ordRank\_const}); the $\Lam$-adic half ($\Mord$ free of finite rank over $\Lam$ and the control
isomorphism for all $s$) is only stated, as the open hypothesis \lean{Langlandscape.Hyp.MordLambdaFree} (Lean-side status: open, S04-L14b). Thus
(a) as a whole stays \stproved{} (tex), with the Lean part \emph{formalized modulo \pathtt{Hyp.MordLambdaFree}} for the $\Lam$-adic half. The
Lean hypothesis \lean{Langlandscape.Hyp.MordLambdaFree} is thus the statement of a tex-proved theorem (not a merely cited one), taken as an argument because the data
$G_n=\Delta_{n+1}$, the reduction maps and the diamond representations are not yet in Lean; it has no named consumer in the stage-05/06 plans.
\emph{(b), Lean form (M9.3 sub-check, recorded):} \lean{Langlandscape.Hyp.HidaControl} is schematic, on abstract \emph{modules} $M_k$, $M_2$, $Q_k$ (the tex theorem is stated on
duals, $P^{\rm ad}_k$-quotients); it carries the guard $p\nmid\varphi(M)$ (M9.3) and $U_1(M)$-type levels only; its arguments \texttt{hk} and \texttt{h$\varphi$} are not used in the
statement (audit item O1 not applied).
\end{theorem}
\begin{proof}[Proof of (a)]
Put $R_s:=\cO[\Delta_s]$ and $M_s:=e\,\cO[X_s(U)]$ ($s\ge1$). By Prop.~\ref{prop:s04-freeA} and Lemma~\ref{lem:s04-hida}(d) (over $\cO$),
$M_s$ is free over $R_s$ of a rank independent of $s$; at $s=1$, $\Delta_1=1$ and $R_1=\cO$, so this rank is $r$. Each $R_s$ is local
(as in the proof of Prop.~\ref{prop:s04-freeA}, with $\cO$ in place of $\Zp$), and the kernel of $R_{s+1}\to R_s$ is contained in
the maximal ideal $\m_{s+1}$ (it is generated by the $[h]-1$, $h\in H$, which lie in the augmentation ideal).

\emph{Compatible bases.} Let $b^{(1)}_1,\dots,b^{(1)}_r$ be an $R_1$-basis of $M_1$. Given an $R_s$-basis $b^{(s)}$ of $M_s$: the
transition map $\pi_*:M_{s+1}\to M_s$ is surjective, since by Lemma~\ref{lem:s04-hida}(d) it factors as
$M_{s+1}\twoheadrightarrow M_{s+1}\otimes_{R_{s+1}}R_s\xrightarrow{\sim}M_s$. Choose lifts $b^{(s+1)}_i\in M_{s+1}$ of the $b^{(s)}_i$.
Their images generate $M_{s+1}\otimes_{R_{s+1}}R_s=M_{s+1}/\ker(R_{s+1}\to R_s)M_{s+1}$, hence generate $M_{s+1}/\m_{s+1}M_{s+1}$, hence
generate $M_{s+1}$ by Nakayama. So $\phi:R_{s+1}^r\to M_{s+1}$, $(c_i)\mapsto\sum c_ib^{(s+1)}_i$, is surjective; composing with an
isomorphism $M_{s+1}\cong R_{s+1}^r$ gives a surjective endomorphism of $R_{s+1}^r$, which is injective (Vasconcelos), so $\phi$ is
bijective. By induction we obtain bases $b^{(s)}$ with $\pi_*b^{(s+1)}_i=b^{(s)}_i$.

\emph{Freeness.} The bases identify the inverse system $(M_s,\pi_*)$ with $(R_s^r,\text{projections})$, compatibly with the
$R_{s+1}$-linear transition maps. Hence $\Mord(U)=\varprojlim_sM_s\cong\varprojlim_sR_s^r=\Lam^r$ as $\Lam$-modules, with basis
$b_i:=(b^{(s)}_i)_s$.

\emph{Control.} The projection $\Mord(U)\to M_s$ is $\Lam$-linear ($\Lam$ acting on $M_s$ through $\Lam\to R_s$) and sends the basis
$b_i$ to the basis $b^{(s)}_i$; so the induced map $\Mord(U)\otimes_\Lam R_s=R_s^r\to M_s$ is an isomorphism. It is the map induced by
the projection, which commutes with $U_p$, $T_\ell$, $S_\ell$, $U_\ell$, $\Pi_\ell$, $\langle d\rangle$: $U_p$ commutes with $\pi_*$ by
Lemma~\ref{lem:s01-Up}(b); $T_\ell,S_\ell,U_\ell,\Pi_\ell$ ($\ell\ne p$) are sums of right translations by elements supported at $\ell$,
which commute with $\pi_*$ since $\pi_*$ only changes the level at $p$; and $\langle d\rangle$ commutes with $\pi_*$ because $\delta_d$
normalises every $UU_s$ (Well-definedness after Def.~\ref{def:s04-Mord}). So the control isomorphism is Hecke equivariant.
\end{proof}
\noindent\emph{Proof of (b):} see Remark~\ref{rem:s04-controlb}.

\subsection{The factor 2}
\noindent\emph{Factor $2$.} At $s\ge1$ there are no stabilisers and $-1\notin UU_s$, so no factor $2$ arises inside the tower.
When comparing with a level containing $-1$ (if $-1\in U^{(p)}$: e.g.\ $U$ itself, $U\cdot\Uzero p_p$, or $\Uzero{q}U$ at $s=0$), the stage~01 rule
applies (Prop.~\ref{prop:s01-adjoint}(c)): for $U'\le U$, $\nu\circ\pi_*=\mathrm{tr}\circ\nu'$ holds \emph{assuming} $-1\in U'$ or
$-1\notin U$; without this assumption, $\mathrm{tr}\circ\nu'=2\,\nu\circ\pi_*$ (\lean{Langlandscape.Brandt.normMap\_pushforward},
\lean{Langlandscape.Brandt.normMap\_pushforward\_two} \stformal{}). Passing from $U\ni-1$ to $UU_s\not\ni-1$ (e.g.\ from $\hat\cO^\times$ or
$\Uzero M$ to $\Uone{Mp^s}$) is exactly the case without the assumption, and the hypothesis ``$-1\in U'$ or $-1\notin U$'' must be
checked at \emph{every} level change used. Inside the tower, $s\ge1$, it holds since $-1\notin UU_{s+1}\subseteq UU_s\not\ni-1$; for
$X'_s\to X_s$ and $X_{s+1}\to X'_s$ likewise. Since $2\in\Zp^\times$ this changes pairings over $\Zp$ only by the unit $2$, but every comparison of pairings
across such a step must carry this factor explicitly. Hsieh's
$\#\Gamma_{N,x}:=\#\big((D^\times\cap x\hat R_N^\times x^{-1})\Q^\times/\Q^\times\big)$ equals our $w_x=|\Gamma_x/\Gamma_x\cap\{\pm1\}|$
\lit{hsieh\_2017\_hida-families-and-p-adic-triple-product-l-functions\_arxiv}{(4.1), p.~29}, whereas the import's \texttt{Nat.card} is
$|\Gamma_x|$, which is $2w_x$ at levels containing $-1$. \stproved{} (given Prop.~\ref{prop:s01-adjoint}(c)).

\begin{remark}[GL$_2$ side, for orientation only]\label{rem:s04-hida86}
The seed's pointer ``[Hi86b] Thm.~3.1, Cor.~3.2--3.3'' is to the \emph{elliptic} theory:
\lit{hida\_1986\_iwasawa-modules-attached-to-congruences-of-cusp-forms}{Thm.~3.1, p.~250 (ordinary Hecke algebras free of finite
rank over $\Lam$, $p\ge5$); Cor.~3.2, p.~250 (control, $k\ge2$ for cusp forms); Cor.~3.3, p.~250 (ideals of old forms)}. Hida's
$P_k$ there is normalised via his action (3.1) of $\Zp^\times$ (p.~248) and differs from both $P_k$ and $P^{\rm ad}_k$ by a shift
and/or $\iota$; translate before use. None of these is a statement about $\XB U$. \stcited.
\end{remark}
